\chapter{Objetivos}
\label{objetivos}

Este capítulo concreta el propósito del proyecto y descompone dicho propósito en objetivos específicos. Su contenido parte de la petición de tema inicial y se adapta al alcance finalmente definido para Nuvlo, centrado en Jira Cloud, la visualización de métricas ágiles y la existencia de una demo \textit{offline} reproducible.

\section{Objetivo general}

El objetivo general del proyecto es construir una aplicación web funcional que permita a equipos ágiles visualizar, analizar y gestionar métricas clave de Scrum y Kanban mediante paneles intuitivos y personalizables. La herramienta busca facilitar una lectura clara del estado del trabajo, apoyar la mejora continua y reducir la dependencia de decisiones basadas únicamente en intuición.

La aplicación se apoya en Jira Cloud como fuente principal de datos reales. Tras la autorización del usuario mediante Atlassian OAuth, el sistema sincroniza información de proyectos, incidencias, \textit{sprints} y cambios de estado, persiste dichos datos en una base de datos propia y calcula métricas como \textit{Velocity}, \textit{WIP}, \textit{Lead Time} y \textit{Cycle Time}. Además, incorpora un modo de demostración \textit{offline} basado en un conjunto de datos local, pensado para validar y presentar la aplicación aunque no haya conexión con Jira.

\section{Objetivos específicos}

Para alcanzar el objetivo general se plantean los siguientes objetivos específicos:

\begin{itemize}
    \item Analizar las necesidades de equipos que trabajan con Scrum o Kanban y que requieren una visión clara de sus métricas de flujo.
    \item Estudiar las limitaciones de herramientas existentes de gestión y analítica ágil, especialmente en relación con coste, personalización, trazabilidad y dependencia de plataformas externas.
    \item Diseñar una arquitectura web separada por capas, con \textit{frontend}, \textit{backend}, base de datos, caché y componentes compartidos que puedan evolucionar de forma ordenada.
    \item Implementar un mecanismo de autenticación basado en Atlassian OAuth 2.0, evitando registro local de usuarios y delegando el acceso en la cuenta Jira del usuario.
    \item Desarrollar un módulo de sincronización con Jira Cloud capaz de importar proyectos, incidencias, \textit{sprints} y transiciones de estado necesarias para el cálculo de métricas.
    \item Persistir los datos sincronizados en PostgreSQL para que la aplicación no dependa de consultar Jira en cada interacción del usuario.
    \item Utilizar Redis como apoyo operativo para cachear respuestas temporales y exponer el estado de sincronización de forma eficiente.
    \item Calcular métricas ágiles relevantes, incluyendo \textit{Velocity}, \textit{WIP}, \textit{Lead Time}, \textit{Cycle Time}, medias y percentiles cuando los datos disponibles lo permitan.
    \item Diseñar una interfaz \textit{responsive} con \textit{dashboard} principal, filtros, \textit{widgets} configurables, vista de tablero, alertas, actividad y configuración.
    \item Implementar reglas de alerta por umbral e historial de actividad para mejorar la trazabilidad del uso de la aplicación.
    \item Incorporar una demo \textit{offline} reproducible mediante CSV o \textit{dataset} local para permitir validación y presentación sin depender de servicios externos.
    \item Aplicar medidas de seguridad coherentes con el uso de OAuth y datos sensibles, incluyendo \textit{cookies} \texttt{httpOnly}, protección CSRF, PKCE, cifrado de \textit{tokens} y exclusión de secretos del repositorio.
    \item Validar la implementación mediante pruebas unitarias, pruebas de integración, pruebas E2E y comprobaciones manuales sobre Jira real y demo \textit{offline}.
    \item Elaborar la documentación técnica y de usuario necesaria para comprender, ejecutar, mantener y evaluar la aplicación.
\end{itemize}

\section{Alcance de los objetivos}

Algunos objetivos previstos inicialmente se han ajustado durante el desarrollo para mantener un alcance realista y verificable. La integración principal se centra en Jira Cloud, mientras que métricas dependientes de repositorios o despliegues, como \textit{Deployment Frequency}, quedan documentadas como posible ampliación futura. Del mismo modo, la validación se apoya en pruebas automatizadas, revisión funcional y datos reales o simulados, sin plantear una implantación prolongada en varios equipos externos.

Este alcance permite concentrar el trabajo en una solución completa y evaluable: autenticación con Atlassian, sincronización controlada, persistencia local, cálculo de métricas, visualización, alertas, registros de actividad, demo \textit{offline} y documentación asociada.
